\section{Walks, paths, and shortest objects}\label{ch07:sec:walks}
How can we travel from one vertex to another by moving along edges? To answer this precisely, we need names for routes through a graph.

A \emph{walk} from a vertex $u$ to a vertex $w$ is a finite sequence of vertices $v_0,v_1,\ldots,v_k$ that starts at $v_0=u$, ends at $v_k=w$, and in which each two consecutive vertices are joined by an edge: $v_{i-1}v_i$ is an edge for every $i=1,\ldots,k$. The \emph{length} of the walk is $k$, the number of steps it takes. This is one less than the number of entries in the list. A walk may visit the same vertex more than once and may use the same edge more than once; each use of an edge counts as a separate step. A \emph{path} is a walk in which no vertex appears twice. A single vertex $v_0$, with no steps at all, counts as a walk and a path of length $0$ from $v_0$ to itself.

In the example graph $G$, the sequence $a,c,b,c,d$ is a walk of length $4$ from $a$ to $d$. Each consecutive pair is an edge: $ac$, $cb$, $bc$, and $cd$. The walk uses the edge $bc$ twice, once from $c$ to $b$ and once back, and it visits $c$ twice, so it is not a path. The sequence $a,c,d$ is a path of length $2$ from $a$ to $d$. The sequence $a,d$ is not a walk at all, because $ad$ is not an edge. Since edges have no direction, reading a walk backward gives a walk in the opposite direction: $d,c,a$ is a walk from $d$ to $a$. We call two vertices \emph{connected} when there is a walk from one to the other.

A walk that repeats a vertex contains a detour that can be cut out. In the walk $a,c,b,c,d$, the vertex $c$ appears in positions $1$ and $3$, counting positions from $0$. Between these two visits the walk makes the round trip $c,b,c$, which ends where it began. Delete this round trip, keeping a single copy of $c$, and what remains is $a,c,d$. The step from $c$ to $d$ used to follow the second visit to $c$; now it follows the first visit. It is still a legitimate step, because both visits are at the same vertex $c$. The shortened walk still runs from $a$ to $d$, and it is two steps shorter. The next lemma uses this operation to show that a walk of smallest possible length can never repeat a vertex.

\par\noindent\begin{minipage}{\linewidth}
\begin{lemma}\label{ch07:lem:shortest-walk}
Let $u$ and $w$ be vertices of a graph, and suppose that there is a walk from $u$ to $w$. Then there is a walk from $u$ to $w$ of smallest possible length, and every such shortest walk is a path. In particular, there is a path from $u$ to $w$.
\end{lemma}
\end{minipage}\par

\begin{proof}
The length of a walk is a natural number. The set of lengths of walks from $u$ to $w$ is nonempty, because at least one such walk exists. By the well-ordering property (Section~\ref{ch04:sec:contrapositive}, extended to natural numbers in Section~\ref{ch05:sec:not-circular}), this set has a smallest member $L$. Now let $v_0,v_1,\ldots,v_L$ be any walk from $u$ to $w$ of this smallest length $L$.

Suppose, for contradiction, that some vertex appears twice in this walk: $v_i=v_j$ for some positions $i<j$. Delete the entries $v_{i+1},\ldots,v_j$, leaving the sequence
\[
v_0,\ \ldots,\ v_i,\ v_{j+1},\ \ldots,\ v_L.
\]
If $j=L$, nothing comes after $v_i$, and the sequence simply ends at $v_i$. We check that this new sequence is a walk from $u$ to $w$.
\begin{itemize}
\item Each pair of consecutive entries is joined by an edge. Every such pair was already consecutive in the original walk, except for the new pair $v_i,v_{j+1}$, which occurs only when $j<L$. Since $v_i=v_j$, the edge $v_iv_{j+1}$ is the same as $v_jv_{j+1}$, which is an edge of the original walk.
\item The sequence starts at $v_0=u$.
\item It ends at $w$. If $j<L$, its last entry is $v_L=w$. If $j=L$, its last entry is $v_i=v_j=v_L=w$.
\end{itemize}
Its length is $L-(j-i)$, because $j-i$ entries were deleted. Since $j>i$, this is less than $L$, which contradicts the choice of $L$ as the smallest length. Hence no vertex appears twice. Since the shortest walk was arbitrary, every shortest walk from $u$ to $w$ is a path.
\end{proof}

When $u$ and $w$ are connected, the \emph{distance} from $u$ to $w$ is the length of a shortest walk from $u$ to $w$. By Lemma~\ref{ch07:lem:shortest-walk}, this is also the length of a shortest path from $u$ to $w$: a shortest walk is itself a path, and no path can be shorter than the shortest walk, because every path is a walk. In the example graph, the distance from $a$ to $d$ is $2$. The path $a,c,d$ has length $2$, and no shorter walk exists: a walk of length $0$ would need $a=d$, and a walk of length $1$ would need $ad$ to be an edge. Notice that establishing this distance took two facts, a walk of length $2$ and a reason why no shorter walk exists.

The proof of Lemma~\ref{ch07:lem:shortest-walk} never built a path step by step. It chose a walk that was as short as possible and showed that this choice already has the property we wanted, because any repeated vertex would allow a cut that produces an even shorter walk. An object chosen to be as small or as large as possible in some measured respect is called \emph{extremal}, and an argument that starts from such a choice is called an \emph{extremal argument}. We have met this idea before: Chapter~\ref{ch:proof} chose a fraction for $\sqrt2$ with the smallest possible denominator, and Chapter~\ref{ch:induction} chose a smallest counterexample. The next two sections give two more examples, a longest path and a vertex with as many outgoing edges as possible.
